\documentclass[12pt,a4paper]{article}
\usepackage[utf8]{inputenc}
\usepackage[T1]{fontenc}
\usepackage[brazil]{babel}
\usepackage[margin=2.5cm]{geometry}
\usepackage{hyperref}
\usepackage{xcolor}
\usepackage{listings}
\usepackage{booktabs}
\usepackage{enumitem}

\lstset{
  basicstyle=\ttfamily\small,
  backgroundcolor=\color{gray!10},
  frame=single, rulecolor=\color{gray!40},
  columns=fullflexible, breaklines=true,
  literate={á}{{\'a}}1 {ã}{{\~a}}1 {é}{{\'e}}1 {ç}{{\c{c}}}1 {ó}{{\'o}}1
}

\title{Sistema Web de Recomendação para Formação de Equipes\\
\large Repositório e instruções de execução local}
\author{Eduardo Requia Souza \\ Orientador: Alexandre de Oliveira Zamberlan \\
Sistemas de Informação --- UFN}
\date{\today}

\begin{document}
\maketitle

\section{Repositório}

Código-fonte completo (frontend, backend, motor de predição e scripts de avaliação):

\begin{center}
\url{https://github.com/Edurequia/tcc-sistema-recomendacao-equipes}
\end{center}

\begin{tabular}{@{}ll@{}}
\toprule
Pasta & Conteúdo \\
\midrule
\texttt{frontend/} & Interface web (React + TypeScript) \\
\texttt{backend/} & Orquestrador (Java 21, Spring Boot 3, Spring AI) \\
\texttt{api\_motor\_predicao/} & Motor de predição (Python 3.10, FastAPI, PyCaret/LightGBM) \\
\texttt{avaliacao/} & Carga de dados e execução dos quatro cenários \\
\bottomrule
\end{tabular}

\medskip
O sistema não possui autenticação: ao abrir o endereço local, todas as telas ficam acessíveis, sem usuário e senha.

\section{Pré-requisitos}

\begin{itemize}[noitemsep]
  \item Git;
  \item \textbf{Python 3.10} (exigido pelo PyCaret 3.3.2);
  \item \textbf{Java 21} (JDK). O Maven não precisa ser instalado: o projeto usa o Maven Wrapper;
  \item \textbf{Node.js 20} ou superior;
  \item Opcional: chave gratuita da API do Gemini (\url{https://aistudio.google.com/apikey}), necessária apenas para a justificativa textual e o assistente conversacional.
\end{itemize}

Não é necessário instalar banco de dados: as instruções abaixo usam o perfil \texttt{dev}, com banco H2 em memória.

\section{Passo a passo}

São necessários três terminais, um para cada serviço. Os comandos abaixo são para Windows (PowerShell). No Linux ou macOS, troque \texttt{.\textbackslash.venv\textbackslash Scripts\textbackslash python.exe} por \texttt{.venv/bin/python} e \texttt{.\textbackslash mvnw} por \texttt{./mvnw}.

\subsection*{1. Clonar o repositório}
\begin{lstlisting}
git clone https://github.com/Edurequia/tcc-sistema-recomendacao-equipes.git
cd tcc-sistema-recomendacao-equipes
\end{lstlisting}

\subsection*{2. Motor de predição (terminal 1)}
O modelo treinado (\texttt{motor\_predicao\_tcc.pkl}) já está na pasta.
\begin{lstlisting}
cd api_motor_predicao
python -m venv .venv
.\.venv\Scripts\python.exe -m pip install -r requirements.txt
.\.venv\Scripts\python.exe -m uvicorn app.main:app --port 8000
\end{lstlisting}
Verificação: \url{http://localhost:8000/health} deve responder \texttt{"modeloCarregado": true}.

\subsection*{3. Backend (terminal 2)}
\begin{lstlisting}
cd backend
.\mvnw spring-boot:run "-Dspring-boot.run.profiles=dev"
\end{lstlisting}
A primeira execução baixa as dependências e leva alguns minutos. Verificação: \url{http://localhost:8080/actuator/health} deve responder \texttt{UP}. A documentação da API fica em \url{http://localhost:8080/swagger-ui.html}.

\paragraph{Opcional: ativar o Gemini.} Antes de iniciar o backend, crie o arquivo \texttt{backend/.env} com:
\begin{lstlisting}
SPRING_AI_MODEL_CHAT=google-genai
GEMINI_API_KEY=sua-chave-aqui
\end{lstlisting}
Sem a chave, o sistema funciona normalmente; apenas a justificativa e o assistente informam que o serviço não está configurado.

\subsection*{4. Dados de exemplo (terminal 3)}
O banco começa apenas com o catálogo de tecnologias. Para cadastrar desenvolvedores e projetos fictícios:
\begin{lstlisting}
cd avaliacao
python popular_base.py
\end{lstlisting}
O script usa apenas a biblioteca padrão do Python.

\subsection*{5. Frontend (terminal 3)}
\begin{lstlisting}
cd ../frontend
npm install
npm run dev
\end{lstlisting}
Abra \url{http://localhost:5173} no navegador.

\section{Roteiro sugerido de teste}

\begin{enumerate}[noitemsep]
  \item Em \textbf{Projetos}, abra um projeto cadastrado e gere a recomendação;
  \item Observe o ranking dos candidatos, a equipe proposta, as métricas da equipe e a comparação com a linha de base aleatória;
  \item Confirme a equipe: as horas dos integrantes passam a ser descontadas nos próximos projetos;
  \item Com o Gemini ativo, gere a justificativa da equipe ou descreva um projeto em linguagem natural no \textbf{Assistente}.
\end{enumerate}

\section{Reproduzir os cenários de avaliação}

Os resultados das quatro situações controladas apresentadas no texto podem ser regenerados. Com o motor de predição no ar e o backend \textbf{recém-iniciado} no perfil \texttt{dev} (base vazia, sem executar o passo 4):
\begin{lstlisting}
cd avaliacao
python cenarios.py
\end{lstlisting}
As tabelas são gravadas em \texttt{avaliacao/resultados/}. Para repetir, reinicie o backend, o que zera a base.

\section{Problemas comuns}

\begin{itemize}[noitemsep]
  \item \textbf{Erro 503 ao gerar recomendação:} o motor de predição (passo 2) não está em execução;
  \item \textbf{``Não foi possível contatar o servidor'' no frontend:} o backend (passo 3) não está em execução;
  \item \textbf{Falha na instalação do PyCaret:} confira se o Python é a versão 3.10;
  \item \textbf{Não atualize o LightGBM:} a versão 4.7.0 encerra o processo ao prever com o modelo; o \texttt{requirements.txt} fixa a 4.6.0.
\end{itemize}

\end{document}
